%% GENERATED by tools/gen_latex.py -- edit content/agents06.py instead.
%% Build:  latexmk -pdf agents06.tex     (or: pdflatex agents06.tex, twice)
\documentclass[aspectratio=169,11pt,t]{beamer}
\usepackage{itbpro}

\coursetitle{Rencana dan refleksi}
\coursesubtitle{Memecah tugas sebelum mengerjakannya, dan memperbaiki hasil sesudahnya --- beserta harga keduanya, yang bisa dihitung.}
\coursekicker{Chapter 6}
\coursesource{Grootendorst \& Alammar, An Illustrated Guide to AI Agents --- bab 6}
\courseduration{3 jam (2 sesi)}
\coursepresenter{Hendri Karisma}{Instructor}
\coursebrand{AI for Professional \textperiodcentered\ ITB \textperiodcentered\ Ch 6}

\title{Rencana dan refleksi}
\author{Hendri Karisma}
\date{}

\begin{document}
\covertitle

\framekicker{Learning outcomes}
\begin{frame}[shrink=25]{Session Objectives}
\leadpar{By the end of this session, participants can:}
\begin{isteps}
\item \textbf{Membedakan tiga tingkat perencanaan}, dan menyebutkan kapan yang paling sederhana sudah cukup.
\item \textbf{Menghitung biaya satu gelung refleksi}, dan menyebutkan syarat henti yang harus dipasang sebelum menyalakannya.
\item \textbf{Membedakan sinyal dari dunia dan sinyal dari model}, dan menjelaskan kenapa hanya yang pertama menghasilkan perbaikan terarah.
\item \textbf{Menyebutkan kapan menyusun ulang rencana}, dan kenapa itu perlu anggarannya sendiri.
\item \textbf{Menyebutkan tiga cara refleksi memperburuk hasil}.
\item \textbf{Merancang perbaikan berkelanjutan yang aman}: apa yang boleh berubah sendiri dan apa yang harus lewat manusia.
\end{isteps}

\vskip 6pt
\vskip 4pt
\reslink{Site}{Course home}{https://hendrikarisma.my.id/rs/ai-products-course/index.html}
\reslink{Repo}{ai-agentic-demo}{https://github.com/situkangsayur/ai-agentic-demo}
\reslink{Book}{Maarten Grootendorst \& Jay Alammar, An Illustrated Guide to AI Agents (O'Reilly Media, Early Release)}{https://www.oreilly.com/library/view/an-illustrated-guide/9798341662681/}
\vskip 2pt
\end{frame}

\coursesection{01}{Merencanakan}{Memecah tugas sebelum mengerjakannya --- dan kapan itu justru merugikan.}

\framekicker{Rencana}
\begin{frame}[shrink=25]{Gelung tanpa rencana tidak tahu kapan ia hampir selesai}
\begin{center}
\includegraphics[width=0.92\linewidth,height=0.60\textheight,keepaspectratio]{figs/agents06-planfirst.tex.pdf}
\end{center}
\vskip 3pt{\color{ink3}\fontsize{7.4}{9.4}\selectfont Sama-sama berlangkah; hanya satu yang punya gagasan tentang akhirnya.\par}
{\fontsize{9.4}{12.6}\selectfont\color{fg2}Gelung dasar memilih satu langkah pada satu waktu. Itu cukup untuk banyak tugas, dan punya satu kelemahan yang khas: \textbf{tidak ada apa pun di dalamnya yang tahu berapa banyak pekerjaan yang tersisa.} Ia berhenti karena tujuannya tercapai, atau karena anggarannya habis.\par}\vskip 4pt
\begin{iband}[signal]
Rencana eksplisit menambahkan satu hal sederhana: \textbf{daftar langkah yang bisa ditandai selesai.} Dari situ muncul kemajuan yang bisa dilaporkan, dan syarat henti yang bukan sekadar anggaran.
\end{iband}
\end{frame}

\framekicker{Rencana}
\begin{frame}[shrink=25]{Tiga tingkat, dan yang pertama sering sudah cukup}
\begin{center}
\includegraphics[width=0.92\linewidth,height=0.60\textheight,keepaspectratio]{figs/agents06-levels.tex.pdf}
\end{center}
\vskip 3pt{\color{ink3}\fontsize{7.4}{9.4}\selectfont Naik satu tingkat hanya kalau tingkat sebelumnya terbukti gagal.\par}
\begin{itable}{>{\raggedright\arraybackslash}p{0.2256\textwidth}>{\raggedright\arraybackslash}p{0.376\textwidth}>{\raggedright\arraybackslash}p{0.3384\textwidth}}
\thead{Tingkat} & \thead{Cocok kalau} & \thead{Harganya} \\ \midrule
Tanpa rencana & Langkahnya sedikit dan alatnya jelas & Tidak ada tambahan \\
Rencana di awal & Langkahnya banyak, urutannya penting & Satu panggilan model tambahan \\
Rencana + susun ulang & Amatan sering mengubah arah & Satu panggilan tiap kali disusun ulang \\
\end{itable}
\begin{iband}[signal]
Kesalahan yang paling sering: memasang tingkat ketiga sejak awal karena terdengar canggih. \textbf{Naiklah satu tingkat hanya setelah jejak menunjukkan tingkat sebelumnya gagal}, dan gagal dengan cara yang rencana bisa perbaiki.
\end{iband}
\end{frame}

\framekicker{Rencana}
\begin{frame}[shrink=25]{Rencana yang berguna punya tiga sifat}
\begin{minipage}[t]{0.3253\textwidth}
\begin{icardgood}
\cardh{Bisa ditandai selesai}
Tiap langkah punya kriteria yang jelas. \textbackslash\{\}u201cPahami masalahnya\textbackslash\{\}u201d tidak bisa ditandai; \textbackslash\{\}u201cambil mutasi 12 bulan\textbackslash\{\}u201d bisa.
\end{icardgood}
\end{minipage}%
\hfill
\begin{minipage}[t]{0.3253\textwidth}
\begin{icardgood}
\cardh{Menyebut alat}
Langkah yang tidak bisa dipetakan ke alat mana pun adalah langkah yang tidak bisa dikerjakan.
\end{icardgood}
\end{minipage}%
\hfill
\begin{minipage}[t]{0.3253\textwidth}
\begin{icardgood}
\cardh{Sedikit}
Rencana dua belas langkah biasanya berarti tugasnya belum dipahami, bukan berarti tugasnya besar.
\end{icardgood}
\end{minipage}%
{\fontsize{9.4}{12.6}\selectfont\color{fg2}Kartu kedua yang paling sering dilanggar. Rencana yang berisi langkah tanpa alat menghasilkan agen yang \textbf{mengarang kemajuan}: ia menandai langkah selesai padahal tidak ada yang terjadi.\par}\vskip 4pt
\end{frame}

\framekicker{Rencana}
\begin{frame}[shrink=25]{Menyusun ulang, dan kenapa itu butuh anggarannya sendiri}
\begin{center}
\includegraphics[width=0.92\linewidth,height=0.60\textheight,keepaspectratio]{figs/agents06-replan.tex.pdf}
\end{center}
\vskip 3pt{\color{ink3}\fontsize{7.4}{9.4}\selectfont Amatan boleh mengubah rencana --- tetapi tidak tanpa batas.\par}
{\fontsize{9.4}{12.6}\selectfont\color{fg2}Rencana yang disusun sebelum apa pun diketahui pasti akan salah pada sebagian tugas. Membiarkannya diperbarui saat amatan masuk adalah perbaikan yang nyata.\par}\vskip 4pt
\begin{iband}[signal]
Dan ia menciptakan gelung baru yang bisa berputar: \textbf{susun rencana, kerjakan satu langkah, susun ulang, kerjakan satu langkah\ldots{}} Anggaran penyusunan ulang harus dipasang terpisah dari anggaran langkah, atau agen bisa sibuk merencanakan tanpa mengerjakan.
\end{iband}
\end{frame}

\framekicker{Rencana}
\begin{frame}[shrink=25]{Ketika rencana justru merugikan}
\begin{minipage}[t]{0.494\textwidth}
\begin{icardbad}
\cardh{Tugas satu atau dua langkah}
Menyusun rencana untuk sesuatu yang selesai dalam satu panggilan alat adalah biaya murni.
\end{icardbad}
\end{minipage}%
\hfill
\begin{minipage}[t]{0.494\textwidth}
\begin{icardbad}
\cardh{Ketika hasil langkah pertama menentukan segalanya}
Rencana yang disusun sebelum melihat data akan salah, dan kesalahannya \textbackslash\{\}u201cresmi\textbackslash\{\}u201d --- agen cenderung mengikutinya.
\end{icardbad}
\end{minipage}%


\vskip 4pt

\begin{minipage}[t]{0.494\textwidth}
\begin{icardbad}
\cardh{Rencana jadi belenggu}
Agen menandai langkah selesai supaya rencananya habis, bukan supaya tugasnya beres.
\end{icardbad}
\end{minipage}%
\hfill
\begin{minipage}[t]{0.494\textwidth}
\begin{icard}
\cardh{Konteks bertambah}
Rencana ikut di konteks tiap giliran, dan tiap versi barunya menambah lagi.
\end{icard}
\end{minipage}%
\begin{iband}[signal]
Kartu ketiga bentuk kegagalan yang halus dan nyata: \textbf{rencana berubah dari alat bantu jadi definisi keberhasilan.} Terlihat di jejak sebagai langkah yang ditandai selesai tanpa panggilan alat apa pun.
\end{iband}
\end{frame}

\coursesection{01b}{Memecah dan mengurutkan}{Dua pekerjaan berbeda yang sering dianggap satu.}

\framekicker{Memecah}
\begin{frame}[shrink=25]{Uji yang menentukan sebuah pecahan sudah cukup kecil}
\begin{center}
\includegraphics[width=0.92\linewidth,height=0.60\textheight,keepaspectratio]{figs/agents06-decomp.tex.pdf}
\end{center}
\vskip 3pt{\color{ink3}\fontsize{7.4}{9.4}\selectfont Pecah sampai tiap bagian bisa dipetakan ke satu alat.\par}
{\fontsize{9.4}{12.6}\selectfont\color{fg2}Memecah tugas terdengar seperti keterampilan yang kabur. Sebenarnya ia punya uji yang tegas: \textbf{sebuah bagian sudah cukup kecil kalau ia bisa dipetakan ke satu alat dan satu hasil yang bisa diperiksa.} Dan kalau masih ada bagian yang tidak bisa dipetakan ke alat mana pun, yang kurang bukan pemecahannya --- \hilite{alatnya yang kurang}.\par}\vskip 4pt
\end{frame}

\framekicker{Memecah}
\begin{frame}[shrink=25]{Dua arah memecah, dan yang salah menghasilkan rencana yang tak terpakai}
\begin{minipage}[t]{0.485\textwidth}
{\fontsize{9.4}{12.6}\selectfont\color{fg2}\textbf{Menurut kegiatan} --- ``analisis, lalu evaluasi, lalu simpulkan''.\par}\vskip 4pt
{\fontsize{9.4}{12.6}\selectfont\color{fg2}Terdengar rapi, dan tiap langkahnya tidak bisa ditandai selesai maupun dipetakan ke alat. Ini rencana yang akan dipatuhi secara formalitas.\par}\vskip 4pt

\end{minipage}%
\hfill
\begin{minipage}[t]{0.485\textwidth}
{\fontsize{9.4}{12.6}\selectfont\color{fg2}\textbf{Menurut hasil antara} --- ``dapatkan mutasi 12 bulan; hitung rasio; ambil klausul yang berlaku''.\par}\vskip 4pt
{\fontsize{9.4}{12.6}\selectfont\color{fg2}Tiap langkah punya keluaran yang konkret, punya alatnya, dan gagalnya kelihatan.\par}\vskip 4pt

\end{minipage}%
\begin{iband}[signal]
Cara cepat memeriksa rencana yang dihasilkan agen: \textbf{bisakah tiap langkah gagal dengan cara yang bisa disebutkan?} Kalau tidak, langkah itu tidak melakukan apa-apa.
\end{iband}
\end{frame}

\framekicker{Mengurutkan}
\begin{frame}[shrink=25]{Pikir, tindak, amati --- dan kenapa urutannya penting}
\begin{center}
\includegraphics[width=0.92\linewidth,height=0.60\textheight,keepaspectratio]{figs/agents06-react.tex.pdf}
\end{center}
\vskip 3pt{\color{ink3}\fontsize{7.4}{9.4}\selectfont Menuliskan alasan sebelum bertindak membuat alasannya jadi bagian dari perhitungan.\par}
{\fontsize{9.4}{12.6}\selectfont\color{fg2}Model menuliskan apa yang kurang, memanggil alat, lalu membaca hasilnya. Alasannya bukan estetika --- \textbf{teks yang ditulis jadi konteks untuk langkah berikutnya}, jadi alasan yang ditulis lebih dulu ikut menentukan pilihannya. Karena itu meminta penjelasan \hilite{sesudah} alat dipanggil menghasilkan pembenaran, bukan alasan.\par}\vskip 4pt
\end{frame}

\framekicker{Mengurutkan}
\begin{frame}[shrink=25]{Bentuk jejak yang membuat sisanya mungkin}
{\fontsize{9.4}{12.6}\selectfont\color{fg2}Kalau tiap giliran mencatat empat hal yang sama, hampir semua yang dibahas di bab 7 jadi bisa dikerjakan tanpa tambahan apa pun.\par}\vskip 4pt
\begin{itable}{>{\raggedright\arraybackslash}p{0.2068\textwidth}>{\raggedright\arraybackslash}p{0.376\textwidth}>{\raggedright\arraybackslash}p{0.3572\textwidth}}
\thead{Dicatat} & \thead{Dipakai untuk} & \thead{Kalau tidak ada} \\ \midrule
Alasan singkat & Menelusuri kenapa alat itu dipilih & Hanya bisa menebak \\
Panggilan + argumen & Menghitung ulang, menguji, mengaudit & Tidak ada bukti apa yang terjadi \\
Hasil (atau pengenalnya) & Memeriksa angka di jawaban & Angka tanpa asal-usul \\
Anggaran terpakai & Melihat bentuk proses, bukan hanya hasilnya & Biaya jadi kejutan bulanan \\
\end{itable}
\begin{iband}[signal]
Empat kolom ini yang membuat demo bisa menjawab pertanyaan ``atas dasar apa'' setahun kemudian --- dan itu satu-satunya bentuk penjelasan yang bertahan di depan pemeriksa.
\end{iband}
\end{frame}

\framekicker{Mengurutkan}
\begin{frame}[shrink=25]{Tiga cara pola ini dipasang, dari yang paling murah}
\begin{itable}{>{\raggedright\arraybackslash}p{0.2256\textwidth}>{\raggedright\arraybackslash}p{0.3572\textwidth}>{\raggedright\arraybackslash}p{0.3572\textwidth}}
\thead{Cara} & \thead{Yang dibutuhkan} & \thead{Kapan pantas} \\ \midrule
Lewat prompt & Tidak ada & Hampir selalu --- model modern sudah dilatih untuk ini \\
Penyetelan terbimbing & Ribuan jejak yang benar & Kalau domainnya sangat khusus dan polanya sering meleset \\
Penguatan & Pemeriksa otomatis + infrastruktur & Jarang, dan hampir tidak pernah di luar penyedia model \\
\end{itable}
{\fontsize{9.4}{12.6}\selectfont\color{fg2}Perhatikan kolom terakhir. Untuk hampir semua tim, pertanyaannya bukan \textbf{cara mana} --- melainkan apakah deskripsi alat dan bentuk jejaknya sudah cukup baik, sebab keduanya memperbaiki hal yang sama dengan biaya yang jauh lebih kecil.\par}\vskip 4pt
\end{frame}

\framekicker{Rencana}
\begin{frame}[shrink=25]{Rencana yang ditunjukkan adalah kontrak dengan penggunanya}
{\fontsize{9.4}{12.6}\selectfont\color{fg2}Rencana punya kegunaan yang sering terlewat karena bukan soal teknis: ia bisa \textbf{ditunjukkan sebelum dikerjakan}. Pengguna melihat empat langkah, menyetujui atau mengoreksi, lalu agen berjalan.\par}\vskip 4pt
\begin{minipage}[t]{0.485\textwidth}
{\fontsize{9.4}{12.6}\selectfont\color{fg2}\textbf{Menyetujui rencana}\par}\vskip 4pt
{\fontsize{9.4}{12.6}\selectfont\color{fg2}Satu keputusan, di depan, ketika biayanya masih nol. Pengguna melihat maksud sistem sebelum ada efek.\par}\vskip 4pt

\end{minipage}%
\hfill
\begin{minipage}[t]{0.485\textwidth}
{\fontsize{9.4}{12.6}\selectfont\color{fg2}\textbf{Menyetujui tiap langkah}\par}\vskip 4pt
{\fontsize{9.4}{12.6}\selectfont\color{fg2}Enam keputusan, tiap kali menunggu. Bab 1 sudah menyebut akibatnya: persetujuan yang terlalu sering berubah jadi klik otomatis.\par}\vskip 4pt

\end{minipage}%
\begin{iband}[signal]
Untuk tugas yang punya efek, ini sering \textbf{bentuk pengawasan yang paling bisa dijalankan}: satu tinjauan yang benar-benar dibaca, bukan enam yang tidak.
\end{iband}
\end{frame}

\framekicker{Rencana}
\begin{frame}[shrink=25]{Kenapa langkah yang lebih sedikit hampir selalu menang}
{\fontsize{9.4}{12.6}\selectfont\color{fg2}Bab 1 menghitungnya: sembilan langkah yang masing-masing benar 95\% menghasilkan sistem yang benar sekitar \textbf{63\%}. Rencana yang panjang bukan tanda ketelitian --- ia perkalian peluang gagal.\par}\vskip 4pt
\begin{isteps}
\item \{'h': 'Gabungkan langkah yang selalu berurutan', 'p': 'Kalau dua langkah tidak pernah dipakai terpisah, keduanya satu alat yang belum ditulis.'\}
\item \{'h': 'Pindahkan yang pasti ke kode', 'p': 'Aritmetika, penyaringan, pemformatan. Tidak ada langkah rencana yang pantas dipakai untuk hal yang deterministik.'\}
\item \{'h': 'Periksa di tengah, bukan hanya di ujung', 'p': 'Satu pemeriksaan setelah langkah ketiga memotong rantai kesalahan sebelum ia berlipat.'\}
\end{isteps}
\begin{iband}[signal]
Ukuran yang berguna dan jarang dipakai: \textbf{rerata langkah per tugas yang berhasil.} Kalau ia naik dari waktu ke waktu tanpa perubahan kode, ada yang bergeser di hulu.
\end{iband}
\end{frame}

\framekicker{Rencana}
\begin{frame}[shrink=25]{Rencana yang ditulis manusia mengalahkan rencana yang disusun model}
{\fontsize{9.4}{12.6}\selectfont\color{fg2}Untuk tugas yang \textbf{berulang dan bentuknya tetap}, menyusun rencana dengan model tiap kali adalah biaya berulang untuk jawaban yang sudah diketahui.\par}\vskip 4pt
\begin{itable}{>{\raggedright\arraybackslash}p{0.282\textwidth}>{\raggedright\arraybackslash}p{0.2444\textwidth}>{\raggedright\arraybackslash}p{0.4136\textwidth}}
\thead{Bentuk tugas} & \thead{Rencananya} & \thead{Alasannya} \\ \midrule
Berulang, langkahnya tetap & \textbf{Ditulis, jadi kode} & Deterministik, gratis, bisa diuji seperti kode biasa \\
Berulang, urutannya bervariasi & Cetakan + isian & Kerangka tetap, bagian yang berubah diisi model \\
Benar-benar baru tiap kali & Disusun model & Di sinilah biayanya sepadan \\
\end{itable}
\begin{iband}[signal]
Baris pertama itu tempat sebagian besar pekerjaan bisnis berada, dan ia mengembalikan kita ke Bab 1: \textbf{kalau langkahnya sudah diketahui, yang Anda butuhkan alur tetap, bukan agen.}
\end{iband}
\end{frame}

\framekicker{Rencana}
\begin{frame}[shrink=25]{Langkah yang tidak saling bergantung boleh berjalan bersamaan}
{\fontsize{9.4}{12.6}\selectfont\color{fg2}Rencana membuat satu hal jadi terlihat yang tidak terlihat pada gelung langkah-demi-langkah: \textbf{langkah mana yang tidak butuh hasil langkah lain.}\par}\vskip 4pt
\begin{minipage}[t]{0.485\textwidth}
{\fontsize{9.4}{12.6}\selectfont\color{fg2}\textbf{Berurutan wajib}\par}\vskip 4pt
{\fontsize{9.4}{12.6}\selectfont\color{fg2}Hitung rasio $\rightarrow$ ambil kebijakan yang sesuai rasio. Yang kedua butuh yang pertama.\par}\vskip 4pt

\end{minipage}%
\hfill
\begin{minipage}[t]{0.485\textwidth}
{\fontsize{9.4}{12.6}\selectfont\color{fg2}\textbf{Boleh bersamaan}\par}\vskip 4pt
{\fontsize{9.4}{12.6}\selectfont\color{fg2}Ambil data pengajuan, ambil mutasi rekening, ambil riwayat kredit. Tidak ada yang menunggu.\par}\vskip 4pt

\end{minipage}%
{\fontsize{9.4}{12.6}\selectfont\color{fg2}Menjalankan yang kedua bersamaan memotong waktu dinding tanpa mengurangi token --- biayanya sama, jawabannya sama, dan pengalamannya jauh berbeda.\par}\vskip 4pt
\begin{iband}[signal]
Syaratnya satu, dan Bab 8 kembali ke sana: \textbf{hanya untuk alat baca.} Menjalankan alat tulis bersamaan menciptakan pertanyaan urutan yang tidak ada jawabannya.
\end{iband}
\end{frame}

\coursesection{02}{Refleksi}{Memeriksa hasil sendiri --- dan berapa panggilan yang dibayar untuk itu.}

\framekicker{Refleksi}
\begin{frame}[shrink=25]{Bentuk dasarnya, dan harganya}
\begin{center}
\includegraphics[width=0.92\linewidth,height=0.60\textheight,keepaspectratio]{figs/agents06-reflect.tex.pdf}
\end{center}
\vskip 3pt{\color{ink3}\fontsize{7.4}{9.4}\selectfont Draf, kritik, revisi --- dan pertanyaan yang menentukan: cukup?\par}
\begin{itable}{>{\raggedright\arraybackslash}p{0.3196\textwidth}>{\raggedright\arraybackslash}p{0.3102\textwidth}>{\raggedright\arraybackslash}p{0.3102\textwidth}}
\thead{Putaran refleksi} & \thead{Panggilan model} & \thead{Biaya relatif} \\ \midrule
0 --- langsung jawab & 1 & 1$\times$ \\
1 & 3 & \textbf{3$\times$} \\
2 & 5 & 5$\times$ \\
3 & 7 & 7$\times$ \\
\end{itable}
{\fontsize{9.4}{12.6}\selectfont\color{fg2}Tiap putaran menambah dua panggilan: satu untuk mengkritik, satu untuk merevisi. Dan tiap panggilan membawa seluruh riwayat yang sudah bertambah --- jadi biayanya naik lebih cepat daripada kolom ketiga.\par}\vskip 4pt
\end{frame}

\framekicker{Refleksi}
\begin{frame}[shrink=25]{Sinyalnya menentukan apakah ini perbaikan atau perubahan}
\begin{center}
\includegraphics[width=0.92\linewidth,height=0.60\textheight,keepaspectratio]{figs/agents06-signal.tex.pdf}
\end{center}
\vskip 3pt{\color{ink3}\fontsize{7.4}{9.4}\selectfont Dua sumber kritik, dan hanya satu yang tahu sesuatu.\par}
\begin{minipage}[t]{0.485\textwidth}
{\fontsize{9.4}{12.6}\selectfont\color{fg2}\textbf{Sinyal dari dunia} --- uji yang gagal, skema yang ditolak, klausul kebijakan yang dilanggar, hitungan yang tidak cocok.\par}\vskip 4pt
{\fontsize{9.4}{12.6}\selectfont\color{fg2}Ini memberi tahu \textbf{apa} yang salah, jadi revisinya terarah. Di sinilah refleksi benar-benar bekerja.\par}\vskip 4pt

\end{minipage}%
\hfill
\begin{minipage}[t]{0.485\textwidth}
{\fontsize{9.4}{12.6}\selectfont\color{fg2}\textbf{Sinyal dari model sendiri} --- model diminta menilai jawabannya sendiri.\par}\vskip 4pt
{\fontsize{9.4}{12.6}\selectfont\color{fg2}Model yang tidak tahu jawabannya salah saat menulis juga tidak tahu saat memeriksa. Yang dihasilkan sering \textbf{perubahan}, bukan perbaikan.\par}\vskip 4pt

\end{minipage}%
\begin{iband}[signal]
Aturan yang menghemat banyak uang: \textbf{nyalakan refleksi kalau ada pemeriksa; pikir dua kali kalau kritiknya datang dari model yang sama.}
\end{iband}
\end{frame}

\framekicker{Refleksi}
\begin{frame}[shrink=25]{Tiga cara refleksi memperburuk}
\begin{minipage}[t]{0.3253\textwidth}
\begin{icardbad}
\cardh{Berputar tanpa henti}
Kritik selalu bisa menemukan sesuatu. Tanpa syarat henti, gelungnya hanya berhenti karena anggaran.
\end{icardbad}
\end{minipage}%
\hfill
\begin{minipage}[t]{0.3253\textwidth}
\begin{icardbad}
\cardh{Merevisi yang sudah benar}
Jawaban benar diubah jadi salah karena kritik menuntut perubahan. Ini nyata dan sering.
\end{icardbad}
\end{minipage}%
\hfill
\begin{minipage}[t]{0.3253\textwidth}
\begin{icardbad}
\cardh{Yakin yang bertambah, bukan benar}
Jawaban yang sudah melewati tiga revisi terbaca jauh lebih meyakinkan, tanpa jaminan lebih tepat.
\end{icardbad}
\end{minipage}%
\begin{iband}[signal]
Kartu kedua yang paling merugikan dan paling mudah diukur: \textbf{bandingkan hasil dengan dan tanpa refleksi pada kumpulan uji yang sama.} Kalau ada kasus yang tadinya benar lalu jadi salah, refleksinya sedang merusak.
\end{iband}
\end{frame}

\framekicker{Refleksi}
\begin{frame}[shrink=25]{Syarat henti yang harus dipasang sebelum menyalakannya}
\begin{isteps}
\item \{'h': 'Batas putaran, selalu', 'p': 'Dua biasanya cukup. Tiga jarang menambah apa pun yang bisa diukur.'\}
\item \{'h': 'Berhenti kalau pemeriksa sudah lulus', 'p': 'Kalau kriterianya objektif, lulus berarti selesai --- tidak perlu bertanya apakah masih bisa lebih baik.'\}
\item \{'h': 'Berhenti kalau revisinya tidak berubah banyak', 'p': 'Dua revisi yang nyaris sama berarti gelungnya sudah selesai meskipun kritiknya belum kehabisan kata.'\}
\item \{'h': 'Simpan versi terbaik, bukan versi terakhir', 'p': 'Kalau ada pemeriksa yang memberi skor, ambil yang tertinggi. Versi terakhir belum tentu versi terbaik.'\}
\end{isteps}
\begin{iband}[signal]
Langkah keempat sering terlewat dan gratis: tanpa itu, satu revisi buruk di akhir membuang seluruh kerja tiga putaran sebelumnya.
\end{iband}
\end{frame}

\framekicker{Refleksi}
\begin{frame}[shrink=25]{Memperbaiki draf, atau mengingat kegagalan --- dua hal berbeda}
\begin{minipage}[t]{0.485\textwidth}
{\fontsize{9.4}{12.6}\selectfont\color{fg2}\textbf{Memperbaiki draf}\par}\vskip 4pt
{\fontsize{9.4}{12.6}\selectfont\color{fg2}Kritik dan revisi berlaku pada satu keluaran, lalu dibuang. Tidak ada yang tersisa untuk tugas berikutnya.\par}\vskip 4pt
{\fontsize{9.4}{12.6}\selectfont\color{fg2}Murah dipasang, dan cocok untuk keluaran yang punya pemeriksa.\par}\vskip 4pt

\end{minipage}%
\hfill
\begin{minipage}[t]{0.485\textwidth}
{\fontsize{9.4}{12.6}\selectfont\color{fg2}\textbf{Mengingat kegagalan}\par}\vskip 4pt
{\fontsize{9.4}{12.6}\selectfont\color{fg2}Ketika sebuah percobaan gagal, catatan tentang \textit{kenapa} disimpan dan ikut dibaca pada percobaan berikutnya.\par}\vskip 4pt
{\fontsize{9.4}{12.6}\selectfont\color{fg2}Lebih kuat pada tugas yang boleh dicoba beberapa kali --- dan membawa semua pertanyaan ingatan dari Bab 4.\par}\vskip 4pt

\end{minipage}%
\begin{iband}[signal]
Yang kedua sering disalahpahami sebagai ``agen yang belajar''. Tidak ada bobot yang berubah; yang berubah \textbf{isi konteks pada percobaan berikutnya}, dan itu bisa dihapus.
\end{iband}
\end{frame}

\framekicker{Refleksi}
\begin{frame}[shrink=25]{Kalau ada pemeriksa, refleksi berubah sifat}
{\fontsize{9.4}{12.6}\selectfont\color{fg2}Bab 3 menyebut syaratnya: memeriksa harus lebih mudah daripada menjawab. Ketika syarat itu terpenuhi, gelung refleksi berhenti jadi tebakan dan jadi \textbf{perbaikan terarah}.\par}\vskip 4pt
\begin{isteps}
\item \{'h': 'Jalankan pemeriksanya, bukan tanyakan pendapat model', 'p': 'Uji yang gagal, skema yang ditolak, klausul yang dilanggar --- semuanya menyebut apa yang salah.'\}
\item \{'h': 'Berikan pesan kegagalannya apa adanya', 'p': 'Pesan pemeriksa lebih berguna daripada ringkasan tentangnya.'\}
\item \{'h': 'Berhenti begitu lulus', 'p': 'Kriteria objektif berarti tidak perlu bertanya apakah masih bisa lebih baik.'\}
\end{isteps}
\begin{iband}[signal]
Perhatikan bahwa ini persis bentuk yang sudah dipakai pengembang perangkat lunak: \textbf{jalankan uji, baca kegagalannya, perbaiki, ulangi.} Yang berubah hanya siapa yang menulis perbaikannya.
\end{iband}
\end{frame}

\framekicker{Refleksi}
\begin{frame}[shrink=25]{Mengukur apakah refleksi Anda menghasilkan sesuatu}
\begin{itable}{>{\raggedright\arraybackslash}p{0.2444\textwidth}>{\raggedright\arraybackslash}p{0.3384\textwidth}>{\raggedright\arraybackslash}p{0.3572\textwidth}}
\thead{Bandingkan} & \thead{Yang dilihat} & \thead{Kesimpulan kalau buruk} \\ \midrule
Tanpa vs 1 putaran & Tugas benar, biaya, waktu & Kalau ketepatan tetap, matikan --- 3$\times$ biaya untuk nol \\
Kasus yang berubah arah & Benar $\rightarrow$ salah & Ada berarti refleksinya merusak, bukan sekadar tidak menolong \\
Putaran ke-2 vs ke-1 & Selisihnya & Hampir selalu kecil; batas dua sudah cukup \\
Panjang keluaran & Bertambah tiap putaran? & Revisi yang hanya memanjangkan bukan revisi \\
\end{itable}
{\fontsize{9.4}{12.6}\selectfont\color{fg2}Baris kedua yang paling sering mengejutkan tim yang memasang refleksi karena ``pasti lebih baik'': \textbf{ada kasus yang tadinya benar dan jadi salah}, dan tanpa perbandingan berpasangan itu tidak akan pernah terlihat.\par}\vskip 4pt
\end{frame}

\framekicker{Refleksi}
\begin{frame}[shrink=25]{Kritik dari model yang berbeda, bukan dari model yang sama}
{\fontsize{9.4}{12.6}\selectfont\color{fg2}Kalau kritik memang harus datang dari model, memakai \textbf{model yang berbeda} dari yang menulis lebih masuk akal daripada memakai yang sama.\par}\vskip 4pt
{\fontsize{9.4}{12.6}\selectfont\color{fg2}Alasannya sederhana: kesalahan yang datang dari cara sebuah model memandang masalah cenderung tidak terlihat oleh model itu sendiri, dan lebih mungkin terlihat oleh model yang dilatih berbeda.\par}\vskip 4pt
\begin{iband}[signal]
Tetapi hitung dulu ongkosnya: ini menambah \textbf{satu penyedia lagi} untuk dipantau, dibayar, dan diperbarui. Sering kali pemeriksa berupa kode memberi manfaat yang sama dengan harga jauh lebih murah --- dan hasilnya pasti.
\end{iband}
\end{frame}

\framekicker{Refleksi}
\begin{frame}[shrink=25]{Kapan refleksi jelas sepadan}
\begin{minipage}[t]{0.3253\textwidth}
\begin{icardgood}
\cardh{Keluarannya bisa dijalankan}
Kode, kueri, konfigurasi. Pemeriksanya objektif dan kegagalannya menyebut barisnya.
\end{icardgood}
\end{minipage}%
\hfill
\begin{minipage}[t]{0.3253\textwidth}
\begin{icardgood}
\cardh{Keluarannya harus memenuhi bentuk}
Skema, format laporan, kelengkapan medan. Bisa diperiksa kode, dan revisinya terarah.
\end{icardgood}
\end{minipage}%
\hfill
\begin{minipage}[t]{0.3253\textwidth}
\begin{icardbad}
\cardh{Prosa bebas tanpa kriteria}
Di sini refleksi menghasilkan teks yang berbeda, bukan lebih benar --- dan tiga kali lebih mahal.
\end{icardbad}
\end{minipage}%
{\fontsize{9.4}{12.6}\selectfont\color{fg2}Pola yang muncul lagi: \textbf{manfaat refleksi hampir seluruhnya berasal dari ada atau tidaknya pemeriksa}, bukan dari kecanggihan gelungnya.\par}\vskip 4pt
\end{frame}

\coursesection{03}{Sistem yang memperbaiki diri}{Dan garis yang memisahkan itu dari sistem yang berubah tanpa diawasi.}

\framekicker{Perbaikan diri}
\begin{frame}[shrink=25]{Tiga hal berbeda yang sering disebut dengan satu nama}
\begin{itable}{>{\raggedright\arraybackslash}p{0.282\textwidth}>{\raggedright\arraybackslash}p{0.282\textwidth}>{\raggedright\arraybackslash}p{0.376\textwidth}}
\thead{Yang berubah} & \thead{Umur perubahannya} & \thead{Risikonya} \\ \midrule
Jawaban, dalam satu tugas & Satu proses & Rendah --- refleksi biasa \\
Ingatan / catatan & Antar percakapan & Sedang --- bisa membawa kesalahan ke pengguna lain \\
\textbf{Bobot model} & Permanen & Tinggi --- dan hampir tidak pernah yang dimaksud orang \\
\end{itable}
{\fontsize{9.4}{12.6}\selectfont\color{fg2}Ketika orang berkata \textbackslash\{\}u201cagen yang belajar\textbackslash\{\}u201d, yang hampir selalu dimaksud adalah baris pertama atau kedua. Baris ketiga menuntut jalur pelatihan, data, dan peninjauan --- dan mengubah sistem yang sudah divalidasi.\par}\vskip 4pt
\end{frame}

\framekicker{Perbaikan diri}
\begin{frame}[shrink=25]{Yang boleh berubah sendiri, dan yang tidak}
\begin{minipage}[t]{0.485\textwidth}
{\fontsize{9.4}{12.6}\selectfont\color{fg2}\textbf{Boleh, dengan catatan}\par}\vskip 4pt
\begin{ibullets}
\item Catatan tentang preferensi pengguna
\item Hasil yang mahal dihitung ulang, dengan kedaluwarsa
\item Urutan alat yang terbukti lebih cepat
\end{ibullets}

\end{minipage}%
\hfill
\begin{minipage}[t]{0.485\textwidth}
{\fontsize{9.4}{12.6}\selectfont\color{fg2}\textbf{Lewat manusia, selalu}\par}\vskip 4pt
\begin{ibullets}
\item Alat baru atau izin baru
\item Aturan kebijakan
\item Apa pun yang mengubah batas kemampuan
\end{ibullets}

\end{minipage}%
\begin{iband}[signal]
Garisnya sama dengan Bab 1 dan tidak bergerak: \textbf{sistem boleh memperbaiki caranya bekerja; ia tidak boleh memperluas apa yang boleh dilakukannya.} Yang kedua adalah keputusan orang.
\end{iband}
\end{frame}

\framekicker{Perbaikan diri}
\begin{frame}[shrink=25]{Perbaikan yang paling andal tidak melibatkan model sama sekali}
{\fontsize{9.4}{12.6}\selectfont\color{fg2}Pola yang jauh lebih sering berhasil daripada sistem yang melatih dirinya: \textbf{kumpulkan jejak, temukan pola kegagalan, perbaiki alat dan deskripsinya, jalankan lagi kumpulan uji.}\par}\vskip 4pt
\begin{isteps}
\item \{'h': 'Jejak menunjukkan langkah mana yang paling sering gagal', 'p': 'Biasanya satu atau dua alat menyumbang sebagian besar kegagalan.'\}
\item \{'h': 'Perbaikannya biasanya teks, bukan bobot', 'p': 'Deskripsi alat, pesan galat, urutan konteks. Semuanya bisa ditinjau dan dibatalkan.'\}
\item \{'h': 'Kumpulan uji mengatakan apakah berhasil', 'p': 'Dan menyimpannya sebagai kasus baru mencegah kegagalan yang sama kembali.'\}
\end{isteps}
\begin{iband}[signal]
Ini membosankan, dan itu kelebihannya: \textbf{tiap perubahannya bisa dibaca, ditinjau, dan dikembalikan} --- tiga sifat yang tidak dimiliki sistem yang memperbarui bobotnya sendiri.
\end{iband}
\end{frame}

\framekicker{Perbaikan diri}
\begin{frame}[shrink=25]{Melatih diri dari soal yang dibuatnya sendiri}
{\fontsize{9.4}{12.6}\selectfont\color{fg2}Arah penelitian yang menarik dan pantas dipahami sebelum dipercaya: sistem membuat soal untuk dirinya sendiri, mencoba menjawabnya, memeriksa jawabannya, dan belajar dari yang benar --- tanpa data berlabel dari manusia.\par}\vskip 4pt
\begin{minipage}[t]{0.485\textwidth}
{\fontsize{9.4}{12.6}\selectfont\color{fg2}\textbf{Kenapa bisa bekerja}\par}\vskip 4pt
{\fontsize{9.4}{12.6}\selectfont\color{fg2}Pada domain yang jawabannya bisa diperiksa mesin --- matematika, kode --- memeriksa memang lebih mudah daripada menjawab, jadi imbalannya bisa dihitung tanpa manusia.\par}\vskip 4pt

\end{minipage}%
\hfill
\begin{minipage}[t]{0.485\textwidth}
{\fontsize{9.4}{12.6}\selectfont\color{fg2}\textbf{Kenapa belum tentu berlaku pada Anda}\par}\vskip 4pt
{\fontsize{9.4}{12.6}\selectfont\color{fg2}Soal yang dibuat sistem cenderung mirip yang sudah bisa dijawabnya, dan tanpa pemeriksa objektif seluruh gelungnya kehilangan jangkarnya.\par}\vskip 4pt

\end{minipage}%
\begin{iband}[signal]
Pertanyaan yang memisahkan yang bisa dipakai dari yang belum: \textbf{apakah domain Anda punya pemeriksa yang tidak melibatkan model?} Kalau tidak, sistem yang melatih dirinya sedang menilai dirinya dengan alat ukur buatannya sendiri.
\end{iband}
\end{frame}

\framekicker{Praktik}
\begin{frame}[shrink=25]{Memasang rencana dan refleksi tanpa merusak yang sudah jalan}
\begin{isteps}
\item \{'h': 'Pasang jejaknya dulu', 'p': 'Empat kolom tadi. Tanpa itu, tidak ada cara mengetahui apakah rencana atau refleksi menolong.'\}
\item \{'h': 'Tambahkan satu, ukur, baru tambahkan yang lain', 'p': 'Menyalakan keduanya sekaligus membuat perubahan angka tidak bisa diatribusikan.'\}
\item \{'h': 'Batasi sejak menit pertama', 'p': 'Anggaran langkah, anggaran susun-ulang, batas putaran refleksi. Ketiganya lebih mudah dipasang sekarang daripada sesudah tagihan pertama.'\}
\item \{'h': 'Simpan kasus yang berubah arah', 'p': 'Kasus yang jadi salah setelah refleksi adalah kasus uji terbaik yang akan Anda dapatkan.'\}
\end{isteps}
\end{frame}

\framekicker{Praktik}
\begin{frame}[shrink=25]{Empat bentuk kegagalan, dan tandanya di jejak}
\begin{itable}{>{\raggedright\arraybackslash}p{0.2444\textwidth}>{\raggedright\arraybackslash}p{0.3572\textwidth}>{\raggedright\arraybackslash}p{0.3384\textwidth}}
\thead{Kegagalan} & \thead{Tandanya di jejak} & \thead{Perbaikannya} \\ \midrule
Merencanakan tanpa mengerjakan & Banyak susun-ulang, sedikit panggilan alat & Anggaran susun-ulang terpisah \\
Menandai selesai tanpa bukti & Langkah selesai tanpa panggilan alat di antaranya & Kriteria selesai harus menyebut alat \\
Refleksi berputar & Panjang keluaran naik, isinya berputar & Batas putaran, dan simpan versi terbaik \\
Rencana jadi belenggu & Amatan jelas bertentangan, rencana tetap diikuti & Izinkan susun ulang, dengan anggarannya \\
\end{itable}
{\fontsize{9.4}{12.6}\selectfont\color{fg2}Keempatnya terlihat di kolom yang sama --- \textbf{jejak per giliran} --- dan tidak satu pun menimbulkan galat. Pola yang sama seperti bab-bab sebelumnya.\par}\vskip 4pt
\end{frame}

\framekicker{Praktik}
\begin{frame}[shrink=25]{Kapan tidak usah memasang keduanya}
\begin{minipage}[t]{0.494\textwidth}
\begin{icardgood}
\cardh{Tugasnya pendek dan jelas}
Dua sampai tiga panggilan alat, urutannya sudah tertentu. Rencana dan refleksi hanya menambah biaya dan tempat gagal.
\end{icardgood}
\end{minipage}%
\hfill
\begin{minipage}[t]{0.494\textwidth}
\begin{icardgood}
\cardh{Sudah ada pemeriksa di jalur utama}
Kalau validasi dan aturan kebijakan sudah menolak yang salah, sebagian besar manfaat refleksi sudah didapat lebih murah.
\end{icardgood}
\end{minipage}%
{\fontsize{9.4}{12.6}\selectfont\color{fg2}Demo kredit UMKM tidak memakai rencana eksplisit maupun gelung refleksi. Urutannya tetap, tiap langkah punya alatnya, dan pemeriksaannya kode --- jadi keduanya akan menambah biaya tanpa menambah apa pun yang bisa diukur.\par}\vskip 4pt
\begin{iband}[signal]
Ini kesimpulan yang berulang di seluruh modul, dan pantas diulang: \textbf{mekanisme yang lebih canggih hampir selalu kalah oleh alat yang dirancang lebih baik.}
\end{iband}
\end{frame}

\framekicker{Praktik}
\begin{frame}[shrink=25]{Menilai mutu rencana tanpa menilai jawabannya}
{\fontsize{9.4}{12.6}\selectfont\color{fg2}Rencana bisa dinilai sendiri, terpisah dari hasil akhirnya --- dan itu berguna, sebab rencana yang buruk kadang tetap menghasilkan jawaban benar secara kebetulan.\par}\vskip 4pt
\begin{itable}{>{\raggedright\arraybackslash}p{0.3196\textwidth}>{\raggedright\arraybackslash}p{0.282\textwidth}>{\raggedright\arraybackslash}p{0.3384\textwidth}}
\thead{Pertanyaan} & \thead{Diperiksa dengan} & \thead{Kalau gagal} \\ \midrule
Tiap langkah punya alat? & Kode --- cocokkan dengan daftar alat & Rencananya mengandung langkah kosong \\
Tiap langkah bisa ditandai selesai? & Kode --- ada kriteria? & Akan ditandai selesai tanpa bukti \\
Langkahnya sedikit? & Hitung & Peluang gagalnya berlipat \\
Urutannya mungkin? & Kode --- periksa ketergantungan & Langkah butuh hasil yang belum ada \\
\end{itable}
\begin{iband}[signal]
Keempatnya \textbf{pemeriksaan kode}, bukan penilaian model --- jadi keempatnya bisa jadi uji yang berjalan tiap commit, dan tidak butuh jaringan.
\end{iband}
\end{frame}

\framekicker{Praktik}
\begin{frame}[shrink=25]{Rencana dan refleksi menambah dua anggaran baru}
{\fontsize{9.4}{12.6}\selectfont\color{fg2}Bab 1 menyebut enam cara gelung harus bisa berhenti. Memasang rencana dan refleksi menambahkan dua lagi, dan keduanya sering lupa dipasang justru karena terasa seperti fitur, bukan risiko.\par}\vskip 4pt
\begin{minipage}[t]{0.494\textwidth}
\begin{icardaccent}
\cardh{Anggaran susun-ulang}
Berapa kali rencana boleh diganti dalam satu proses. Tanpa ini, agen bisa sibuk merencanakan tanpa pernah mengerjakan.
\end{icardaccent}
\end{minipage}%
\hfill
\begin{minipage}[t]{0.494\textwidth}
\begin{icardaccent}
\cardh{Anggaran putaran refleksi}
Berapa kali kritik-revisi boleh berjalan. Kritik selalu bisa menemukan sesuatu, jadi tanpa batas ia tidak akan berhenti sendiri.
\end{icardaccent}
\end{minipage}%
\begin{iband}[signal]
Keduanya harus \textbf{terpisah} dari anggaran langkah. Kalau digabung, proses yang boros merencanakan akan kehabisan langkah sebelum mengerjakan apa pun --- dan gejalanya terbaca seperti tugasnya terlalu sulit.
\end{iband}
\end{frame}

\framekicker{Praktik}
\begin{frame}[shrink=25]{Urutan yang menghemat waktu, sekali lagi dari belakang}
\begin{isteps}
\item \{'h': 'Alat yang baik dulu', 'p': 'Sebagian besar tugas yang ``butuh perencanaan'' sebenarnya butuh satu alat yang mengerjakan tiga langkah sekaligus.'\}
\item \{'h': 'Jejak per giliran', 'p': 'Tanpa itu, semua yang di bab ini hanya bisa dinilai dari perasaan.'\}
\item \{'h': 'Pemeriksa, kalau domainnya punya', 'p': 'Ini yang menentukan apakah refleksi akan berguna atau hanya mahal.'\}
\item \{'h': 'Baru rencana, baru refleksi', 'p': 'Satu per satu, dengan pengukuran di antaranya.'\}
\end{isteps}
\begin{iband}[signal]
Kalau urutan ini diikuti, sebagian tim menemukan bahwa mereka \textbf{tidak pernah sampai ke langkah keempat} --- dan sistemnya sudah cukup baik. Itu hasil yang benar, bukan pekerjaan yang belum selesai.
\end{iband}
\end{frame}

\framekicker{Penutup}
\begin{frame}[shrink=25]{Yang dibawa pulang dari bab ini}
\begin{isteps}
\item \{'h': 'Rencana memberi gagasan tentang akhir', 'p': 'Tapi naik tingkat hanya setelah jejak menunjukkan tingkat sebelumnya gagal.'\}
\item \{'h': 'Satu putaran refleksi berharga tiga panggilan', 'p': 'Dan tiap panggilan membawa riwayat yang lebih panjang, jadi biayanya naik lebih cepat daripada hitungan panggilannya.'\}
\item \{'h': 'Refleksi bekerja kalau kritiknya datang dari dunia', 'p': 'Model yang tidak tahu jawabannya salah saat menulis juga tidak tahu saat memeriksa.'\}
\item \{'h': 'Simpan versi terbaik, bukan versi terakhir', 'p': 'Gratis, dan menyelamatkan kerja beberapa putaran.'\}
\item \{'h': 'Sistem boleh memperbaiki caranya, bukan memperluas izinnya', 'p': 'Garis yang sama sejak Bab 1, dan ia tidak bergerak.'\}
\end{isteps}
\end{frame}
\note{Latihan: ambil satu tugas di sistem mereka, jalankan dengan dan tanpa satu putaran refleksi, dan hitung tiga angka. Sering hasilnya mengejutkan yang mengusulkannya.}

\end{document}
